\ifpdf
    \graphicspath{{Sections/5_discussion_figs/Raster/}{Sections/5_discussion_figs/PDF/}{Sections/5_discussion_figs/}}
\else
    \graphicspath{{Sections/5_discussion_figs/Vector/}{Sections/5_discussion_figs/}}
\fi

\chapter{Discussion}
\label{sec:discussion}

\section{Comparison with Hu et al.}

We retained PaDIS training and patch-score construction while moving from AAPM parallel-beam to LIDC-IDRI fan-beam data \cite{}. Numerical identity was neither expected nor useful. We instead test which conclusions transfer across dataset and forward operator.

\subsection{Findings that Transferred}

We reproduced several of the main findings. Hu et al. reported PaDIS PSNRs of \(33.57\,\mathrm{dB}\) at 20 views and \(29.48\,\mathrm{dB}\) at 8 views \cite{}. Our paper-style track reached \(32.71\) and \(27.15\,\mathrm{dB}\). The 20-view result is therefore within \(0.86\,\mathrm{dB}\), despite the changed data and geometry, whilst the 8-view result still improves substantially over FDK and our CP-TV baseline. We also found that PaDIS can generate coherent whole CT slices after training only on patches. In comparison, fixed stitching and averaging produced much less credible global anatomy, supporting the use of randomly shifted partitions.

We also reproduced the practical high-resolution argument \cite{} by training and applying a \(512\times512\) PaDIS prior without a whole-image model at this resolution, obtaining \(40.08\,\mathrm{dB}\) PSNR and \(0.960\) SSIM. At 256 resolution, its logged wall time was \(24\%\) below whole-image VE-DPS after half the training time, although concurrent jobs mean this is not an isolated speed benchmark. PaDIS is therefore still useful where whole-image memory or training cost is prohibitive.

\subsection{Findings that Did Not Transfer}

We did not reproduce the central claim that PaDIS consistently outperforms whole-image diffusion \cite{}. In the paper-style track, whole-image diffusion exceeded PaDIS-DPS by \(1.26\,\mathrm{dB}\) at 20 views and \(1.95\,\mathrm{dB}\) at 8 views. The LION-physics extension showed the same reversal under VE-DPS. The additional examples in \autoref{fig:results-ct20-additional} and \autoref{fig:results-ct8-additional} also show that PaDIS is strong, but is not consistently sharper or more faithful than the whole-image prior.

However, this does not mean that the whole-image prior is always better. Within the LION-physics extension, the separately tuned PaDIS VE-DDNM configuration achieved the best 20- and 60-view PSNRs, and beat the whole-image configuration by \(2.30\,\mathrm{dB}\) at 20 views. Langevin and VE-DPS favoured the whole-image model, whilst predictor-corrector changed the ranking depending on the metric. Since settings were validation-selected for each prior and sampler, these comparisons do not isolate the prior architecture, but suggest that PaDIS performance depends on how its patch score is used. We can reproduce the claim that PaDIS is a flexible and sometimes stronger prior, but not that it always outperforms whole-image diffusion.

\subsection{Ablation Findings}

Hu et al. also ablated patch size, training-set size and positional encoding \cite{}. We repeated these on LIDC-IDRI. The initialisation and schedule comparisons instead resolve paper--code differences.

Our patch-size analysis does not show a simple benefit from smaller patches. The \(96\times96\) model had almost the same mean PSNR as whole-image VE-DPS, whilst the default \(56\times56\) model was \(1.43\,\mathrm{dB}\) lower. The non-monotonic improvement at \(8\times8\) suggests that receptive field, optimisation, number of patches and partition geometry all contribute. The results in \autoref{tab:patch-size} and \autoref{fig:results-patch-size} suggest that retaining more global context can be beneficial when memory permits, although the non-default models inherited the \(P=56\) inference setting and may not be at their individual optima. At this image size, the clearest PaDIS advantage is feasibility, rather than improved quality from very small patches.

Using all processed LIDC-IDRI slices did not improve PaDIS. However, we fixed training by wall-clock time, so the full-data model saw each image fewer times. We can only conclude that additional data did not help under equal compute. An equal-epoch comparison may behave differently.

Removing position channels reduced PSNR by only \(0.3\) to \(0.4\,\mathrm{dB}\), unlike the failure reported by Hu et al. \cite{}. The CT measurements and repeated data-consistency updates appear to provide enough global information to locate local anatomy. This finding remains conditional on the shared inference setting and does not show that position is unnecessary for unconditional generation, where no measurements provide location. Similarly, FDK and Gaussian initialisation converged to almost identical results after 1000 model evaluations. The geometric noise schedule was more important, improving PSNR by about \(1.3\,\mathrm{dB}\) over the public EDM-style default.

\subsection{Original PaDIS Specification and Implementation}
\label{sec:discussion-reproducibility}

We treated the released code as part of the specification of the original PaDIS work, rather than as a separate implementation \cite{}. Even together, the original PaDIS paper and code leave several choices ambiguous. The variable-size U-Net is attributed to different references, whilst the number of sampled patches, validation split, checkpoint rule and 25 test slices are not given. View and detector counts are reported without the physical field of view, detector spacing or complete angular conventions needed to define \(A\). The released 180-view fan-beam configuration also spans only \(180^\circ\), which is unclear from the main description. Fine-tuning is reported without the search or selection protocol. There are smaller reporting inconsistencies. Table 1 does not always bold the best SSIM, and Table 7 is captioned as a dataset-size analysis despite comparing samplers.

The released code resolves some details, but differs from the algorithms in the original PaDIS paper \cite{}. Its main DPS path starts from filtered back-projection rather than the Gaussian state in Algorithm 1, differentiates the residual norm rather than its square, and defaults to an EDM-style schedule. Predictor-corrector changes the corrector noise level and reuses its patch layout. VE-DDNM uses 100 noise levels with 10 inner updates rather than 1000 separate levels. We also restored several comparison APIs before the published methods could run. These ambiguities within the combined specification can materially change the reconstruction.

We found agreement within sampler variation when comparing against the adapted public-code fork on the same LION geometry. This agreement validates the Public-compatible implementation, but not the unspecified ODL geometry in the original PaDIS paper. Publishing geometry, data indices, configurations and checkpoints would reduce this ambiguity.

\section{Extensions to the Original Work}

We use Paper-style and Public-compatible as two readings of the combined specification, while LION-physics is an extension introduced here that replaces geometry-specific calibration with an operator-derived scale and whose results we use to assess transfer and stabilisation, not direct reproduction of Hu et al.'s algorithms.

\subsection{Operator-Normalised Data Consistency}

The LION-physics DPS rows reached relative sinogram residuals near \(10^{-3}\), over an order of magnitude below many paper-style and public-compatible results. This is important because diffusion models can otherwise produce plausible images which are not well supported by the measurements. In the limited-angle experiment, whole-image VE-DPS gave the best PSNR and SSIM, whilst PaDIS-DPS gave the smallest residual. There is therefore a trade-off between similarity to this particular held-out target and strict agreement with the information which was measured.

The public-compatible fan-beam updates required calibrated gradient scales of \(0.0405\) and \(0.1022\). Replacing these with \(\lVert F\rVert_2^2\) normalisation still gave high-quality results across each angular setting. These constants therefore correct for a particular implementation, rather than defining PaDIS itself. We believe operator normalisation is a more reproducible choice because it remains meaningful when the geometry changes.

\subsection{Cost of Fixed Patch Assembly}

Fixed averaging obtained the strongest displayed SSIM and MAE in parts of \autoref{tab:main-quality}, but these rows contain only four images and have wide uncertainty. They also required roughly 24 minutes per slice, about twelve times the PaDIS-DPS runtime, and their unconditional generations lacked credible global anatomy. They show that overlap handling can suppress seams, but do not provide a practical improvement over random PaDIS partitions.

\section{Limitations and Future Work}

As in the original CT experiments of Hu et al. \cite{}, our sinograms omit photon noise, scatter, beam hardening and motion. We therefore investigate angular undersampling, not realistic low-dose CT. In this work, we use two-dimensional experiments, one dataset and one training run per model, and evaluate 25 slices rather than independent scans. We use four images for fixed-overlap and native-resolution rows. Bootstrap intervals measure variation across slices, not retraining uncertainty, some high-view methods inherit 20-view tuning, and we did not tune the patch-size and no-position rows per model. PSNR, SSIM and MAE also cannot determine whether a nodule was preserved or hallucinated.

Future work should add calibrated photon noise and task-level evaluation of clinically relevant features. Equal-epoch, multi-seed training would separate dataset size from optimisation time. Averaging scores over several partitions, or a lower-variance deterministic estimator, may help with accelerated samplers. External datasets and expert assessment are required before making clinical claims.

Overall, we reproduce the ability of PaDIS to assemble coherent and high-quality reconstructions from patch scores, including at native resolution. We do not reproduce universal superiority over whole-image diffusion. Instead, our results show where PaDIS is useful, which implementation choices change its ranking, and how operator-aware scaling improves transfer between CT geometries.
